\documentclass[a4paper,11pt]{article}
\usepackage[T2A]{fontenc}
\usepackage[utf8]{inputenc}
\usepackage[russian]{babel}
\usepackage{amsmath,amssymb,mathtools}
\usepackage{helvet}
\renewcommand{\familydefault}{\sfdefault}
\usepackage{icomma}
\usepackage[margin=22mm,top=23mm,bottom=23mm,headheight=23pt]{geometry}
\usepackage{microtype}
\usepackage{xcolor}
\usepackage{tikz}
\usetikzlibrary{arrows.meta,positioning,shapes.geometric,calc}
\usepackage{tabularx,booktabs,array}
\usepackage{enumitem}
\usepackage{fancyhdr}
\usepackage{hyperref}
\usepackage{titlesec}
\usepackage{needspace}
\usepackage{setspace}
\definecolor{accent}{HTML}{087F81}
\definecolor{darkaccent}{HTML}{075C61}
\definecolor{textgray}{HTML}{333D40}
\definecolor{muted}{HTML}{667376}
\hypersetup{colorlinks=true,linkcolor=darkaccent,pdfauthor={Лёвин Артём Александрович},pdftitle={Чек-лист. Встречное движение. Ярослав Гаврилов}}
\titleformat{\section}{\Large\color{darkaccent}}{\thesection.}{0.65em}{}[\vspace{3pt}\titlerule]
\titleformat{\subsection}{\large\color{darkaccent}}{}{0pt}{}
\titlespacing*{\section}{0pt}{15pt}{13pt}
\titlespacing*{\subsection}{0pt}{12pt}{5pt}
\setlist[enumerate]{leftmargin=*,itemsep=8pt,topsep=5pt}
\setlist[itemize]{leftmargin=*,itemsep=4pt,topsep=4pt}
\setlength{\parindent}{0pt}
\setlength{\parskip}{6pt}
\linespread{1.12}
\pagestyle{fancy}
\fancyhf{}
\fancyhead[L]{\small\color{muted}Задачи на встречное движение}
\fancyhead[R]{\small\color{muted}Ярослав Гаврилов}
\renewcommand{\headrulewidth}{0.35pt}
\renewcommand{\headrule}{\hbox to\headwidth{\color{accent}\leaders\hrule height \headrulewidth\hfill}}
\fancyfoot[C]{\small\thepage}
\newcommand{\idea}[1]{\par\noindent{\color{darkaccent}\bfseries Ключевая идея.} #1\par}
\newcommand{\warning}[1]{\par\noindent{\color{darkaccent}\bfseries Проверьте.} #1\par}
\newcommand{\example}[1]{\par\noindent{\color{darkaccent}\bfseries Пример.} #1\par}
\newcommand{\qtask}[2]{\needspace{7\baselineskip}\item \textbf{#1} #2}
\begin{document}
\thispagestyle{empty}
\begin{center}
\vspace*{20mm}
{\large\color{darkaccent}МАТЕМАТИКА\par}
\vspace{13mm}
{\LARGE\color{textgray}Чек-лист\par}
\vspace{12mm}
{\huge\color{accent}\bfseries Задачи на встречное движение\par}
\vspace{5mm}
{\Large\color{darkaccent}Скорость, время, расстояние; задержка старта\par}
\vspace{13mm}
{\large Подготовка к ЕГЭ по профильной математике\par}
\vfill
{\Large Ярослав Гаврилов\par}
\vspace{10mm}
{\large 9 ноября 2026 г.\par}
\vspace{13mm}
{\normalsize Лёвин Артём Александрович\par}
{\normalsize эксклюзивно для Ярослава Гаврилова\par}
\vspace{12mm}
\begin{tikzpicture}[x=1cm,y=1cm]
\draw[accent,line width=.8pt] (-5,0) .. controls (-3.7,.58) and (-2.3,.58) .. (0,0) .. controls (2.3,-.58) and (3.7,-.58) .. (5,0);
\end{tikzpicture}
\end{center}
\clearpage

\section{Основа любой задачи о движении}
Два автомобиля выезжают из городов $A$ и $B$ навстречу друг другу. До встречи каждый проезжает собственную часть дороги. В точке встречи эти части образуют всё расстояние между городами.
\begin{center}
\begin{tikzpicture}[>=Latex,x=1cm,y=1cm]
 \coordinate (A) at (0,0); \coordinate (M) at (6,0); \coordinate (B) at (13,0);
 \draw[line width=1.0pt,accent] (A)--(B);
 \fill[darkaccent] (A) circle (2pt) (M) circle (2pt) (B) circle (2pt);
 \node[below=6pt] at (A) {$A$};\node[below=6pt] at (M) {$M$};\node[below=6pt] at (B) {$B$};
 \draw[->,line width=1pt] (.5,.5)--(5.4,.5) node[midway,above] {$s_A$};
 \draw[<-,line width=1pt] (6.5,.5)--(12.5,.5) node[midway,above] {$s_B$};
 \draw[<->,muted] (0,-.72)--(13,-.72) node[midway,fill=white,inner sep=4pt] {$D$};
\end{tikzpicture}
\end{center}
\[
\boxed{s=vt},\qquad t=\frac{s}{v},\qquad v=\frac{s}{t},\qquad \boxed{s_A+s_B=D}.
\]
Здесь $v$~--- скорость в км/ч, $t$~--- время в часах, $s$ и $D$~--- расстояния в километрах. Одну и ту же единицу времени нужно использовать для обоих автомобилей.

\subsection{Таблица «скорость~--- время~--- расстояние»}
\begin{center}
\renewcommand{\arraystretch}{1.35}
\begin{tabularx}{.94\linewidth}{@{}lccc@{}}
\toprule
Автомобиль & $v$, км/ч & $t$, ч & $s=vt$, км\\
\midrule
Из $A$ в $B$ & $v_A$ & $t_A$ & $v_At_A$\\
Из $B$ в $A$ & $v_B$ & $t_B$ & $v_Bt_B$\\
\bottomrule
\end{tabularx}
\end{center}
\idea{Сначала определите, когда стартовал каждый автомобиль и сколько времени он действительно находился в пути до встречи. Затем заполняйте таблицу и записывайте уравнение.}

\section{Одновременный выезд}
Если автомобили отправились одновременно и встретились через $t$ часов, оба были в пути ровно $t$ часов:
\[
 v_At+v_Bt=D,\qquad (v_A+v_B)t=D,\qquad t=\frac{D}{v_A+v_B}.
\]
Сумма скоростей $v_A+v_B$ называется \emph{скоростью сближения} для встречного движения. После нахождения $t$ можно вычислить расстояние от любого города до встречи.
\example{Между городами 560 км. Навстречу выехали автомобили со скоростями 60 и 80 км/ч. Через сколько часов они встретятся?}
\[
60t+80t=560\quad\Longrightarrow\quad140t=560\quad\Longrightarrow\quad t=4.
\]
Проверка: первый автомобиль прошёл $60\cdot4=240$ км, второй~--- $80\cdot4=320$ км; $240+320=560$ км.
\warning{Складывать скорости и сразу делить на их сумму удобно только при одинаковом времени движения. При задержке выезда время у автомобилей различается.}

\section{Когда неизвестна точка встречи}
Пусть расстояние между городами 405 км. Первый автомобиль движется со скоростью 60 км/ч, второй~--- 75 км/ч. Требуется узнать, на каком расстоянии от города $B$ находится точка встречи.

Обозначим через $x$ километров путь второго автомобиля. Тогда первый проехал $405-x$ километров. Поскольку оба стартовали одновременно, время в пути одинаково:
\[
\frac{405-x}{60}=\frac{x}{75}.
\]
Решаем: $75(405-x)=60x$, откуда $30375=135x$, значит, $x=225$ км.
Проверка: $180/60=3$ ч и $225/75=3$ ч. Путь в сумме составляет 405 км.
\idea{При переменной «расстояние» время находят делением расстояния на скорость. Правильное равенство времени часто короче, чем попытка сразу угадать всю таблицу.}

\section{Выезд с задержкой: главная ловушка}
Пусть первый автомобиль выехал из $A$, а второй отправился из $B$ на $\Delta$ часов позже. Обозначим через $x$ время от выезда \emph{первого} автомобиля до встречи. Тогда:
\[
 t_A=x,\qquad t_B=x-\Delta,\qquad x\geq\Delta.
\]
Равенство суммарного пути имеет вид
\[
 \boxed{v_Ax+v_B(x-\Delta)=D}.
\]
\example{Расстояние равно 645 км. Первый автомобиль выехал со скоростью 70 км/ч; спустя три часа навстречу ему выехал второй со скоростью 75 км/ч. Когда они встретятся, считая от первого выезда?}
\[
70x+75(x-3)=645,\quad 145x=870,\quad \boxed{x=6\text{ ч}}.
\]
Второй находился в пути $6-3=3$ часа. Проверка: $70\cdot6+75\cdot3=420+225=645$ км.
\warning{Число «на три часа позже» вычитается из времени первого автомобиля. Перепутанные $x+3$ и $x-3$ дают неверную модель движения.}

\section{Выбор переменной «расстояние»}
Иногда задача спрашивает, где произошла встреча. Удобно сразу обозначить через $x$ расстояние от города $A$ до точки встречи. Тогда оставшаяся часть маршрута равна $D-x$.

Если второй автомобиль выехал на $\Delta$ часов позже, первый находился в пути дольше на $\Delta$ часов:
\[
 \boxed{\frac{x}{v_A}-\frac{D-x}{v_B}=\Delta},\qquad 0\leq x\leq D.
\]
\example{Между городами 380 км. Из $A$ выехал автомобиль со скоростью 50 км/ч. Через час из $B$ навстречу ему выехал автомобиль со скоростью 60 км/ч. На каком расстоянии от $A$ они встретятся?}
\[
\frac{x}{50}-\frac{380-x}{60}=1.
\]
Умножаем на 300: $6x-5(380-x)=300$, то есть $11x=2200$ и $\boxed{x=200\text{ км}}$.
Проверка: $200/50=4$ часа, $180/60=3$ часа. Разница составляет один час.

\section{Когда ищут неизвестную скорость}
В условии иногда дана точка встречи, известна скорость второго автомобиля и указана задержка его выезда. Тогда удобнее вычислять всё \emph{в обратном порядке}:
\begin{enumerate}[label=\arabic*.]
\item Найдите путь второго автомобиля: $s_B=D-s_A$.
\item Найдите его время движения: $t_B=s_B/v_B$.
\item Учтите задержку: $t_A=t_B+\Delta$.
\item Вычислите скорость первого автомобиля: $v_A=s_A/t_A$.
\end{enumerate}
\example{Между городами 458 км. Первый автомобиль выехал из $A$, второй~--- из $B$ через три часа, со скоростью 64 км/ч. Встреча произошла в 330 км от $A$. Найдите скорость первого автомобиля.}
\[
 s_B=458-330=128\text{ км},\qquad t_B=\frac{128}{64}=2\text{ ч},
\]
\[
 t_A=2+3=5\text{ ч},\qquad \boxed{v_A=\frac{330}{5}=66\text{ км/ч}}.
\]
\warning{Ответ всегда сопоставляйте с вопросом. Если нужно найти скорость, число часов или километров само по себе ответом служить не может.}

\clearpage
\thispagestyle{plain}
\begin{center}
{\Large\color{darkaccent}\bfseries Алгоритм решения задачи на встречное движение}\par
\vspace{10mm}
\end{center}
\tikzset{
 flstart/.style={ellipse,draw=darkaccent,thick,minimum width=5.0cm,minimum height=1.0cm,align=center,font=\small},
 flstep/.style={rounded corners=3pt,rectangle,draw=accent,thick,text width=9.1cm,minimum height=1.0cm,align=center,font=\small,inner sep=5pt},
 fldec/.style={diamond,draw=darkaccent,thick,aspect=2.8,text width=3.3cm,align=center,font=\small,inner sep=2pt},
 flarr/.style={-{Latex[length=3mm]},thick,darkaccent}
}
\begin{center}
\begin{tikzpicture}[node distance=12mm and 12mm]
\node[flstart] (a) {Прочитайте условие и отметьте города,\\скорости, расстояние и моменты старта};
\node[flstep,below=of a] (b) {Нарисуйте отрезок $AB$, отметьте точку встречи $M$;\\запишите $s_A+s_B=D$};
\node[fldec,below=15mm of b] (c) {Есть задержка\\выезда?};
\node[flstep,below left=15mm and 0mm of c,text width=4.5cm] (d) {Нет: $t_A=x$,\\$t_B=x$};
\node[flstep,below right=15mm and 0mm of c,text width=4.5cm] (e) {Да: $t_A=x$,\\$t_B=x-\Delta$};
\node[flstep,below=37mm of c] (f) {Заполните таблицу $v$~--- $t$~--- $s$;\\найдите расстояния по формуле $s=vt$};
\node[flstep,below=of f] (g) {Составьте уравнение по условию, решите его;\\если неизвестно расстояние, используйте $t=s/v$};
\node[flstart,below=of g] (h) {Проверьте расстояния, времена, ограничения\\и единицы ответа};
\draw[flarr] (a)--(b);\draw[flarr] (b)--(c);
\draw[flarr] (c.south west)--node[pos=.5,above left,font=\footnotesize]{нет}(d.north);
\draw[flarr] (c.south east)--node[pos=.5,above right,font=\footnotesize]{да}(e.north);
\coordinate (join) at ([yshift=6mm]f.north);
\draw[thick,darkaccent] (d.south) |- (join);
\draw[thick,darkaccent] (e.south) |- (join);
\fill[darkaccent] (join) circle (1.2pt);
\draw[flarr] (join)--(f.north);
\draw[flarr] (f)--(g);\draw[flarr] (g)--(h);
\end{tikzpicture}
\end{center}
\clearpage

\section{Типичные ошибки и самопроверка}
\begin{itemize}
\item \textbf{Разная длительность движения.} При задержке на два часа запись $t_A=t_B$ неверна. Проверьте, кто начал раньше.
\item \textbf{Неправильно выбранные части пути.} Для встречи $s_A+s_B=D$; если $s_A=x$, то $s_B=D-x$.
\item \textbf{Подмена исходных чисел.} Перепишите скорости из текста до вычислений; похожие записи 70 и 75 ведут к разным ответам.
\item \textbf{Непроверенное время.} Для задержки $\Delta$ модель $t_B=x-\Delta$ требует $x\geq\Delta$.
\item \textbf{Смешение часов и минут.} Полтора часа~--- $1{,}5$ ч; 90 минут~--- тоже $1{,}5$ ч.
\item \textbf{Неверная размерность.} $\text{км/ч}\cdot\text{ч}=\text{км}$, а $\text{км}/(\text{км/ч})=\text{ч}$.
\end{itemize}
\subsection{Контрольные вопросы перед записью ответа}
Укажите, от какого выезда отсчитывается найденное время. Убедитесь, что расстояния положительны и в сумме дают весь путь. Проверьте, что более поздний автомобиль ехал меньше времени. Напишите единицу измерения той величины, которую требуется найти.

\clearpage
\section*{Тренировочный блок: 10 задач}
\addcontentsline{toc}{section}{Тренировочный блок}
\noindent\textcolor{darkaccent}{\bfseries Путь А.} Задачи расположены по нарастающей сложности. Во всех сюжетах скорость автомобилей постоянна, остановок в пути нет, автомобили движутся по одной дороге навстречу друг другу. Для каждого задания полезно построить схему маршрута и заполнить таблицу «скорость~--- время~--- расстояние».
\subsection*{Базовый уровень}
\begin{enumerate}[label=\textbf{\arabic*.},leftmargin=*,itemsep=13pt]
\qtask{Два автомобиля на междугородней трассе.}{
Между городами Берёзовск и Лесной проложена прямая автомобильная дорога длиной 480 км. Утром из Берёзовска в Лесной выехал легковой автомобиль со скоростью 55 км/ч. В ту же минуту из Лесного в противоположном направлении отправился второй автомобиль со скоростью 65 км/ч. Оба водителя продолжают движение без остановок до встречи. Через сколько часов после отправления они встретятся и сколько километров от Берёзовска будет до места встречи?}
\qtask{Где встретятся машины?}{
Между посёлками Северный и Южный 525 км. Из Северного в Южный отправился автомобиль, который всё время едет со скоростью 60 км/ч. Одновременно из Южного навстречу ему начал движение автомобиль со скоростью 80 км/ч. На дороге есть пункт обслуживания, и диспетчеру нужно определить точное место предстоящей встречи автомобилей. Сколько километров будет от Южного до места их встречи?}
\qtask{Восстановите неизвестную скорость.}{
Между городами $A$ и $B$ расстояние составляет 540 км. Из города $A$ выехал первый автомобиль с неизвестной постоянной скоростью. В ту же минуту из города $B$ навстречу ему отправился второй автомобиль со скоростью 70 км/ч. Через три часа после выезда водители встретились. Найдите скорость первого автомобиля в километрах в час.}
\end{enumerate}
\clearpage
\subsection*{Средний уровень}
\begin{enumerate}[label=\textbf{\arabic*.},leftmargin=*,start=4,itemsep=13pt]
\qtask{Встреча при позднем выезде.}{
Расстояние между двумя городами составляет 540 км. Ранним утром из первого города во второй выехал автомобиль со скоростью 60 км/ч. Из второго города навстречу ему должен отправиться другой автомобиль со скоростью 80 км/ч, однако тот начал движение ровно через два часа после первого. Оба ехали непрерывно и встретились в пути. Через сколько часов после отправления первого автомобиля произошла встреча?}
\qtask{Найдите место встречи после задержки.}{
Города $A$ и $B$ соединены дорогой длиной 485 км. Из города $A$ в сторону $B$ выехал автомобиль со скоростью 65 км/ч. Второй автомобиль, находившийся в городе $B$, отправился навстречу ему на один час позже и поддерживал скорость 75 км/ч. Диспетчеру необходимо указать расстояние от города $A$ до точки встречи. Сколько километров проехал первый автомобиль до встречи?}
\qtask{Кто сколько времени находился в пути?}{
Между городами $A$ и $B$ 560 км. Сначала из города $A$ выехал автомобиль со скоростью 50 км/ч. Спустя два часа из города $B$ навстречу ему отправился второй автомобиль, который двигался со скоростью 65 км/ч. Водители встретились до прибытия в пункты назначения. Сколько километров от города $A$ находилось место встречи и сколько часов ехал второй водитель?}
\end{enumerate}
\clearpage
\section*{Тренировочный блок: продолжение}
\subsection*{Уровень выше среднего}
\begin{enumerate}[label=\textbf{\arabic*.},leftmargin=*,start=7,itemsep=15pt]
\qtask{Определите задержку отправления.}{
Между городами $A$ и $B$ расстояние равно 600 км. Из города $A$ первым выехал автомобиль со скоростью 60 км/ч. Позднее из города $B$ ему навстречу отправился второй автомобиль со скоростью 80 км/ч. Известно, что встреча произошла на расстоянии 360 км от города $A$, причём оба водителя ехали без остановок. На сколько часов второй автомобиль отправился позже первого?}
\qtask{Скорость по точке встречи.}{
Между городами $A$ и $B$ 584 км. Первым из города $A$ выехал автомобиль, скорость которого требуется определить. Через два часа из города $B$ навстречу ему отправился второй автомобиль со скоростью 68 км/ч. Встреча произошла в точке, расположенной в 312 км от города $A$. Найдите постоянную скорость первого автомобиля в километрах в час.}
\qtask{Определите расстояние между городами.}{
Две машины движутся навстречу друг другу по дороге, соединяющей города $A$ и $B$. Первая выехала из $A$ со скоростью 64 км/ч. Второй автомобиль начал путь из $B$ на полтора часа позже первого и ехал со скоростью 72 км/ч. По данным системы слежения, автомобили встретились ровно через пять часов после отправления первой машины. Найдите расстояние между городами в километрах.}
\end{enumerate}
\subsection*{Повышенный уровень}
\begin{enumerate}[label=\textbf{\arabic*.},leftmargin=*,start=10,itemsep=15pt]
\qtask{Положение автомобилей в момент второго старта.}{
Расстояние между городами $A$ и $B$ составляет 780 км. В 8:00 из города $A$ в сторону $B$ выехал автомобиль со скоростью 70 км/ч. В 10:00 из города $B$ навстречу ему отправился второй автомобиль со скоростью 90 км/ч. Оба водителя постоянно поддерживали указанные скорости. Какое расстояние разделяло автомобили в 10:00? Через сколько часов после 10:00 они встретились? На каком расстоянии от города $A$ произошла встреча?}
\end{enumerate}
\vfill
\noindent\textcolor{darkaccent}{\bfseries Перед сверкой ответов.} Для каждой задачи укажите, когда началось движение каждого автомобиля; проверьте, какие величины выражаются в часах, а какие~--- в километрах. Если ответ~--- расстояние до точки встречи, уточните, от какого города оно отсчитывается.

\clearpage
\section*{Ответы к тренировочному блоку}
\vspace{4mm}
\renewcommand{\arraystretch}{1.5}
\begin{tabularx}{\linewidth}{@{}>{\centering\arraybackslash}p{14mm}X@{}}
\toprule
№ & Ответ\\
\midrule
1 & 4 часа; 220 км от Берёзовска.\\
2 & 300 км от Южного.\\
3 & 110 км/ч.\\
4 & 5 часов после первого выезда.\\
5 & 260 км от города $A$.\\
6 & 300 км от города $A$; второй водитель ехал 4 часа.\\
7 & Второй автомобиль выехал на 3 часа позже.\\
8 & 52 км/ч.\\
9 & 572 км.\\
10 & В 10:00 между автомобилями было 640 км; затем до встречи прошло 4 часа; место встречи~--- в 420 км от города $A$.\\
\bottomrule
\end{tabularx}
\par\vspace{10mm}
\noindent\textcolor{darkaccent}{\bfseries Итоговая памятка}\par
\vspace{2mm}
\begin{tabularx}{\linewidth}{@{}p{.36\linewidth}X@{}}
\toprule
Что известно & Какой приём использовать\\
\midrule
Один момент отправления & Общая длительность пути $t$ и равенство $(v_A+v_B)t=D$.\\
Выезд с задержкой $\Delta$ & Время первого $x$, время второго $x-\Delta$.\\
Известно место встречи & Части пути $x$ и $D-x$; время равно $s/v$.\\
Неизвестна скорость & Сначала найдите расстояние и время, затем $v=s/t$.\\
Требуется проверить решение & Сложите пройденные расстояния и сравните времена.\\
\bottomrule
\end{tabularx}
\vfill
\begin{center}
{\small\color{muted}Лёвин Артём Александрович\par}
{\small\color{muted}эксклюзивно для Ярослава Гаврилова\par}
\end{center}
\end{document}
